Table~\ref{tab:main} gives the main comparison on held-out targets (five seeds, mean $\pm$ std over seeds), Fig.~\ref{fig:budget} plots success against budget, and Table~\ref{tab:extra} adds the two exploratory annealing variants that we introduced \emph{after} seeing the main results (they are therefore post-hoc and not part of the registered hypotheses). Test statistics below are paired $t$-tests across seeds (from \texttt{rh compare} for systems in the main group and from a logged script for cross-group pairs).

\begin{table*}[t]\centering\caption{Main results: fraction of held-out targets whose design folds uniquely to the target, at oracle budget $K$ (succ\_k$K$). Mean $\pm$ std over five seeds.}\label{tab:main}
\begin{tabular}{lccccc}
\toprule
Method & succ\_k1 $\uparrow$ & succ\_k10 $\uparrow$ & succ\_k100 $\uparrow$ & succ\_k1000 $\uparrow$ & $n$ \\
\midrule
Random search & 0.000 $\pm$ 0.000 & 0.000 $\pm$ 0.000 & 0.011 $\pm$ 0.011 & 0.060 $\pm$ 0.032 & 5 \\
Simulated annealing & 0.000 $\pm$ 0.000 & 0.009 $\pm$ 0.009 & 0.086 $\pm$ 0.036 & 0.558 $\pm$ 0.056 & 5 \\
Contact heuristic & \textbf{0.685 $\pm$ 0.046} & \underline{0.685 $\pm$ 0.046} & \underline{0.685 $\pm$ 0.046} & \underline{0.685 $\pm$ 0.046} & 5 \\
\textbf{Conditional AR designer} & \underline{0.622 $\pm$ 0.065} & \textbf{0.822 $\pm$ 0.038} & \textbf{0.865 $\pm$ 0.030} & \textbf{0.900 $\pm$ 0.028} & 5 \\
\bottomrule
\end{tabular}

\end{table*}

\begin{table*}[t]\centering\caption{Exploratory (post-hoc) annealing variants on the same seeds and test targets. ``Graded'': objective $f_{\log}$; ``heur.\ init'': start from the contact-heuristic sequence.}\label{tab:extra}
\begin{tabular}{lccccc}
\toprule
Method & succ\_k1 $\uparrow$ & succ\_k10 $\uparrow$ & succ\_k100 $\uparrow$ & succ\_k1000 $\uparrow$ & $n$ \\
\midrule
SA from contact heuristic & \textbf{0.685 $\pm$ 0.046} & \textbf{0.702 $\pm$ 0.038} & \textbf{0.731 $\pm$ 0.046} & \textbf{0.867 $\pm$ 0.049} & 5 \\
\bottomrule
\end{tabular}

\vspace{2pt}
\begin{tabular}{lccccc}
\toprule
Method & succ\_k1 $\uparrow$ & succ\_k10 $\uparrow$ & succ\_k100 $\uparrow$ & succ\_k1000 $\uparrow$ & $n$ \\
\midrule
SA graded objective & \underline{0.000 $\pm$ 0.000} & \underline{0.004 $\pm$ 0.006} & \underline{0.287 $\pm$ 0.036} & \underline{0.927 $\pm$ 0.031} & 5 \\
SA graded objective, heuristic init & \textbf{0.685 $\pm$ 0.046} & \textbf{0.722 $\pm$ 0.051} & \textbf{0.820 $\pm$ 0.029} & \textbf{0.982 $\pm$ 0.013} & 5 \\
\bottomrule
\end{tabular}

\end{table*}

\begin{figure}[t]\centering\includegraphics[width=\columnwidth]{budget_curves.pdf}
\caption{Success versus oracle budget $K$ (mean $\pm$ std over five seeds) for all systems. The contact heuristic is budget-independent.}\label{fig:budget}
\end{figure}

\subsection*{H1: zero-oracle designs}
\emph{Registered claim:} at $K=1$ the AR designer beats random sequences and the contact heuristic. \emph{Verdict: not supported.} The AR designer reaches \rhval{main/conditional-ar-designer/hp16/succ_k1/mean:3} against \rhval{main/random-search/hp16/succ_k1/mean:3} for random sequences (paired $p=$ \rhval{cmp/main/random-search/hp16/succ_k1/paired_p}), so it clearly learned something about unseen conformations. But the one-line contact heuristic is better, \rhval{main/contact-heuristic/hp16/succ_k1/mean:3} (AR minus heuristic: \rhval{cmp/main/contact-heuristic/hp16/succ_k1/delta:3}, $p=$ \rhval{cmp/main/contact-heuristic/hp16/succ_k1/paired_p:3}); the difference is not distinguishable from zero with five seeds, but its sign is opposite to the hypothesis. At $K=1$ the learned model's per-sample success rate over all drawn samples is \rhval{main/conditional-ar-designer/hp16/sample_rate/mean:3}, a lower-variance estimate of the same quantity. The test-NLL of the true designing sequences is \rhval{main/conditional-ar-designer/hp16/test_nll/mean:3} nats per residue, against \rhval{abl_components/unconditional-ar/hp16/test_nll/mean:3} for the unconditional model. Because the heuristic already succeeds on most targets, the AR model's value at $K=1$ is largely that it matches a strong hand-made rule without being told it.

\subsection*{H2: equal budget $K=100$}
\emph{Verdict: supported (for the registered SA).} AR succeeds on \rhval{main/conditional-ar-designer/hp16/succ_k100/mean:3} of targets against \rhval{main/simulated-annealing/hp16/succ_k100/mean:3} for SA with the binary-degeneracy objective $f$ (difference \rhval{cmp/main/simulated-annealing/hp16/succ_k100/delta:3}, $p=$ \rhval{cmp/main/simulated-annealing/hp16/succ_k100/paired_p}) and \rhval{main/contact-heuristic/hp16/succ_k100/mean:3} for the heuristic (difference \rhval{cmp/main/contact-heuristic/hp16/succ_k100/delta:3}, $p=$ \rhval{cmp/main/contact-heuristic/hp16/succ_k100/paired_p}). Against the exploratory baselines the margin is smaller or absent: SA with the graded objective is far behind at this budget (\rhval{extra_sa_log/sa-graded-objective/hp16/succ_k100/mean:3}) but SA with graded objective and heuristic initialisation reaches \rhval{extra_sa_log/sa-graded-objective-heuristic-init/hp16/succ_k100/mean:3}, which is below AR by \rhval{analysis_paired/paired-comparison/hp16/ar_vs_salogw_k100_delta/mean:3} with $p=$ \rhval{analysis_paired/paired-comparison/hp16/ar_vs_salogw_k100_p/mean:3}, i.e.\ not distinguishable with five seeds.

\subsection*{H3: does SA catch up at $K=1000$?}
\emph{Registered claim: SA at $K=1000$ is not lower than AR. Verdict for the registered SA: refuted} (SA \rhval{main/simulated-annealing/hp16/succ_k1000/mean:3} versus AR \rhval{main/conditional-ar-designer/hp16/succ_k1000/mean:3}, $p=$ \rhval{cmp/main/simulated-annealing/hp16/succ_k1000/paired_p}). \emph{However, the conclusion reverses with a better baseline objective.} Replacing the binary degeneracy penalty by a graded one lifts cold-start SA to \rhval{extra_sa_log/sa-graded-objective/hp16/succ_k1000/mean:3}, AR minus this SA being \rhval{analysis_paired/paired-comparison/hp16/ar_vs_salog_k1000_delta/mean:3} ($p=$ \rhval{analysis_paired/paired-comparison/hp16/ar_vs_salog_k1000_p/mean:3}; no detectable difference), and graded SA started at the heuristic sequence reaches \rhval{extra_sa_log/sa-graded-objective-heuristic-init/hp16/succ_k1000/mean:3}, which is \emph{better} than AR (difference AR minus SA \rhval{analysis_paired/paired-comparison/hp16/ar_vs_salogw_k1000_delta/mean:3}, $p=$ \rhval{analysis_paired/paired-comparison/hp16/ar_vs_salogw_k1000_p/mean:3}). The registered H3 verdict thus appears to depend on the binary penalty: with it, annealing seems to be trapped among degenerate sequences whose ground-state \emph{energy} equals that of the target (the mean energy gap of the returned design is \rhval{main/simulated-annealing/hp16/gap_k100/mean:3} at $K=100$ for registered SA) but whose ground state is not unique. We did not test an AR model with a larger sampling budget or with learned-model-guided annealing, so we cannot say whether the AR advantage at $K\le 100$ would survive a stronger search or a stronger generator at $K=1000$.

\subsection*{Per-regime breakdown and failure cases}
Table~\ref{tab:regime} splits the budget-10 and budget-100 results by target type. The heuristic is much worse on targets with exactly one designing sequence (\rhval{main/contact-heuristic/hp16/succ_k100_single/mean:3}) than on targets with several (\rhval{main/contact-heuristic/hp16/succ_k100_multi/mean:3}); AR is better on both (\rhval{main/conditional-ar-designer/hp16/succ_k100_single/mean:3} and \rhval{main/conditional-ar-designer/hp16/succ_k100_multi/mean:3}), so its gain over the heuristic is concentrated on the hard, single-solution targets, where it still fails more often (the paired differences within this split were not tested; the number of single-solution targets per seed is about \rhval{main/conditional-ar-designer/hp16/n_single_targets/mean:1}). Failures of the heuristic are almost never wrong-energy designs: its returned design has a ground-state energy equal to the target's in all but a handful of cases (mean energy gap \rhval{main/contact-heuristic/hp16/gap_k1/mean:3}); it fails because the ground state is degenerate. Registered SA shows the same pattern (zero mean energy gap at $K=100$ but low success), consistent with degeneracy rather than wrong folds being its main failure.

\begin{table*}[t]\centering\caption{Success by target type at $K=10$ and $K=100$: ``single'' targets have exactly one uniquely folding sequence, ``multi'' targets several. Mean $\pm$ std over five seeds.}\label{tab:regime}
\begin{tabular}{lccccc}
\toprule
Method & succ\_k100\_single $\uparrow$ & succ\_k100\_multi $\uparrow$ & succ\_k10\_single $\uparrow$ & succ\_k10\_multi $\uparrow$ & $n$ \\
\midrule
Random search & 0.000 $\pm$ 0.000 & 0.017 $\pm$ 0.017 & 0.000 $\pm$ 0.000 & 0.000 $\pm$ 0.000 & 5 \\
Simulated annealing & 0.056 $\pm$ 0.034 & 0.103 $\pm$ 0.054 & 0.000 $\pm$ 0.000 & 0.013 $\pm$ 0.013 & 5 \\
Contact heuristic & \underline{0.522 $\pm$ 0.155} & \underline{0.759 $\pm$ 0.067} & \underline{0.522 $\pm$ 0.155} & \underline{0.759 $\pm$ 0.067} & 5 \\
\textbf{Conditional AR designer} & \textbf{0.819 $\pm$ 0.076} & \textbf{0.886 $\pm$ 0.021} & \textbf{0.742 $\pm$ 0.074} & \textbf{0.859 $\pm$ 0.037} & 5 \\
\bottomrule
\end{tabular}

\end{table*}
